\section{The range \texorpdfstring{\(N\ge e^{15}\)}{N at least exp(15)}}\label{sec:large-tail}

\begin{proposition}\label{prop:large-tail}
Suppose \(U=\log N\ge15\), \(a,d\ge1\), \(d<3a\), and
\[
  s=\frac{d^2}{a+d}<2U+2,\qquad 6a<(5N-11)d .
\]
Then \(R_{N-1}\) and \(R_N\) are coprime.
\end{proposition}

\begin{proof}
Suppose there is a monic irreducible common factor \(F\).
The notation \(m,b,M,h,H,\alpha,V,\Lambda\) is as in
Lemma~\ref{lem:resultant}. Set \(\sigma=2U+2\), \(t=\log(1/\alpha)\),
and \(B=H^4/M\). The starting-point inequality of
Lemma~\ref{lem:degree-cap} is available throughout this range,
and \(\log m>U-.001\).
The shape inequality gives \(b<(5/6)dN\), hence \(t<U\),
\(d<\sigma N\), and \(V<3.24U\).
We first force the full set of rotations in the Padé norm.
The angular gain then rules out small \(H\).
For the remaining factors, counts of square multiples in short
intervals bound the large-prime contribution to the norm.

\subsection*{Rotations and the range of the parameters}

Suppose first that \(q=d/k_0\ge2\) in Lemma~\ref{lem:kummer},
and set \(P=M^{1/(q+1)}\). The right side of
\eqref{eq:pade-norm} is bounded by \(S(X)+\pi(X)V/N\) for any
\(X\ge P\). If \(\alpha\in[\ell,u]\subseteq[.03,.75]\),
the inequality \(M=dN/\alpha<PN/(4\alpha)\) gives
\(P^q<N/(4\alpha)\). Take \(X=(N/(4\ell))^{1/2}\), so
\(X>P\), \(M<X^3\), and \(V/\log X<3.001\).
Using \(S(X)<\log X-.49\) and the prime count in
Lemma~\ref{lem:prime-estimates}, the surplus of the baseline
\(\log m-1.01-g\) over the upper bound is greater than
\[
 \tfrac12(U+\log(4\ell))-1.011-\kappa(2U+2)+.49
                  -\frac{3.6}{\sqrt{4\ell N}},
\]
where \(g\le\kappa s\). Use \(\kappa=.134,.160,7/32\),
respectively, for \([.03,.1]\), \([.1,1/3]\), and
\([1/3,.75]\), by \eqref{eq:pade-loss}.
The smallest endpoint surplus at \(U=15\) is greater than
\(.12\), and the expressions increase with \(U\).

For \(\alpha<.03\) and \(q\ge3\), the inequalities
\(M<dN^2<PN^2/4\) give \(P<X=N^{2/3}/4^{1/3}\)
and \(M<X^4\). The baseline exceeds \(.744U-1.267\).
The upper bound \(\log X-.54+4.6X/N\) gives a surplus
greater than \(.85\) at \(U=15\), increasing thereafter.
If \(\alpha<.03\) and \(q=2\), we have \(P/N<1/2\) and
\(V/\log P<3.001\). The digit loss \(\pi(P)V/N\) is less
than \(4.848P/N<2.43\), since \(P>2H\ge8\).
For \(d<50\) this improves to \(.121\) using
\(P/N<(50/N)^{1/3}\). In that case the gain exceeds \(U-1.22\)
and the prime sum is below
\(\max\{8,(2U+\log50)/3\}-.54\). For \(d\ge50\), the
weaker bound in Lemma~\ref{lem:angles}, with \(n=d/2\),
gives gain greater than \(.9715U+.96t-1.608\) in logarithm.
Indeed \(g_0/n<.049\), \(1/n\le.04\),
\(r/n<.111\), and \(g_0<1.266+.256U\).
The prime sum and digit loss are less than
\(\max\{8,(U+2t+\log\sigma)/3\}+1.89\).
The surplus is positive at \(U=15,t=3.5\), and increases
with both \(U\) and \(t\). We conclude that \(k_0=k=d\).

For this full coset, suppose \(\alpha\ge.03\).
If \(b=2\), then \(a=d=1\), and the norm primes are at
most \(\sqrt{2N}\). The prime-count bound and
\(S(\sqrt{2N})<\log\sqrt{2N}-.49\) already contradict
the baseline \(U-1.011-.176/2\).
Otherwise \(b\ge3\), so \(\log\sqrt M>8\). Since
\(b\le\sigma/(.03)^2<N/4\), the degree and prime bounds
give prime support at most \(\sqrt M\). The right side of
\eqref{eq:pade-norm} is then less than
\begin{equation}\label{eq:large-moderate-upper}
 \tfrac12(U+\log b)-.54+2.35\sqrt{b/N}.
\end{equation}
For \(d\ge10\), at least \(.45\) of the rotations lie in the
left half-plane. On the first and last intervals
\([\ell,u]\) in the table below, the log gain therefore exceeds
\(U-.001-.55(1.01+\kappa\sigma)+.45\log((1-u)/u)\).
On the middle interval at least \(17/30\) of the rotations
have \(|\arg z|\ge\pi/3\), hence \(\Re z\le am\). There the
gain is at least \(me^{-1.01}\), so the logarithmic mean
exceeds \(U-1.011-(13/30)\kappa\sigma\).
Using \(b\le\sigma/\ell^2\), the limiting comparisons are
\[
\begin{array}{c|c|cc}
{}[\ell,u]&\kappa&\text{gain lower bound at }U=15
                    &\text{upper bound at }U=15\\ \hline
{}[.03,.1]&.134&13.07&12.45\\
{}[.1,.5]&.176&11.54&11.08\\
{}[.5,.75]&7/32&10.09&9.41
\end{array}
\]
The affine lower bounds have slope at least \(.75\), whereas
the logarithmic part of the upper bounds has slope at most
\(1/2+1/\sigma\) and the square-root term decreases.
For \(d\le9\),
use \(s<6.75,\ b\le300\) and the baseline
\(U-1.011-7s/32\) against \eqref{eq:large-moderate-upper}.
Hence \(\alpha<.03\) is necessary.

If \(b\le N\), the norm primes are at most \(X=1.03\sqrt M\).
The right side of \eqref{eq:pade-norm} is less than
\(\tfrac12(U+\log b)+1.91\) by the same prime estimates.
When \(d\ge10\), the weaker angular bound gives
\(\log\Lambda>.945U+.9t-1.70\). Here one uses
\(\log C>t-1.417\) and at most \(.211g_0\) in the loss.
Since \(\log b=2t+\log s\le U\), the surplus is greater than
\(.395U-.45\log\sigma-3.61>0\).
If \(2\le d\le9\), the left third and \(s<.27\) give gain
greater than \(U+t/3-.71\) in logarithm, enough since
\(\log b=t+\log d\) and \(t<U\).
For \(d=1\), the last gain in Lemma~\ref{lem:resultant}
exceeds \(U+\tfrac12\log a-1.195\) in logarithm.
Against \(\tfrac12(U+\log b)+1.91\) this leaves a positive
margin, using \(b=a+1\le2a\). The shape inequality precludes
\(b>N\) when \(d=1\), so now \(b>N\) and \(d\ge2\).

For any \(x\ge\max(e^8,\log K)\), the prime sum in
\eqref{eq:pade-norm} is below \(\log x+.461\). Below \(x\)
use Lemma~\ref{lem:prime-estimates}, and above \(x\) bound it
by \(\log K/(x-1)<1.001\).
Suppose \(B<2048\). The equalities
\(B=h^4s^2 b/N=s^3 h^4e^{2t-U}\) imply
\[
 s<8\sqrt2,\qquad V<\tfrac83U+1.62,\qquad
 H<10.08e^{(U+t)/3}.
\]
The last bound follows from \(H=(B/h)^{1/3}e^{(U+t)/3}\).
Split the count of distinct norm primes at \(N/16\)
and use \(\log K\le HV\). The digit loss is less than
\(.236+9.46UH/N\). Both coefficients follow from the prime count at \(U=15\).
To see the decrease, divide the bound \(\tfrac83U+1.62\)
for \(V\) by \(U-\log16\) in the first coefficient, and use
\((\tfrac83+1.62/U)^2/(1-\log16/U)\) in the second.
For \(d\ge50\) the weaker angular bound gives
\(\log\Lambda>U+.98t-1.72\), using \(g_0<2.459\)
and \(g_0/d<.025\).
Splitting the prime sum at \(\max(e^8,HV)\), the surplus is
greater than
\[
 \tfrac23 U+(.98-\tfrac13)t
 -\log(10.08(\tfrac83 U+1.62))-2.417
 -95.3568\,Ue^{(t-2U)/3}.
\]
This is concave in \(t\). At \(t=3.5,U\) its values are
greater than \(3.59,1.60\), respectively, at \(U=15\)
and increase thereafter. If \(d<50\), then \(s<2500/N<.001\)
and \(t>U-\log50\). The gain exceeds \(U+t/3-.71\).
Use \(M<50N^2\), \(H<(102400)^{1/4}N^{1/2}\), and
\(V<2U+\log51\). The digit bound is below \(1.65\), so
the total upper bound is less than
\[
 \tfrac12U+\tfrac14\log102400+\log(2U+\log51)+.461+1.65.
\]
It is beaten at \(U=15\) by the gain, which increases faster.
Thus \(B\ge2048\).

Finally, the initial block at order \(N\) has prime support
below \(1.03b\), since \(b>N\) and \(\alpha<.03\).
Its logarithm exceeds \(H\log(aN)\). Counting valuations
by the affine bound, it is at most
\(H S(1.03b)+\pi(1.03b)\log M\).
The difference \(\log(aN)-S(1.03b)\) exceeds \(U+.47\)
and the prime-count term is below \(2.24b\). We have obtained
\begin{equation}\label{eq:large-parameter-range}
 k_0=k=d,\qquad \alpha<.03,\qquad b>N,\qquad
 B\ge2048,\qquad H/b<\frac{2.24}{U+.47}.
\end{equation}

\subsection*{Norm and degree bounds}

Put \(P=2H+d\) and \(L=\log M\). For primes \(p>P\) one
has \(p^4>M\). Call a prime exceptional if either terminal
interval of length \(6H\) contains a positive multiple of
\(p^3\). For each such interval and each \(q\le M/P^3\)
there is at most one prime with \(qp^3\) in the interval.
Each exceptional prime contributes at most \(L\) to the log
norm by \eqref{eq:norm-affine}. Write \(Y\) for the sum of
\(\log p\) over the nonexceptional unequal-height norm primes.
By Lemma~\ref{lem:two-heights} and \eqref{eq:unequal-products},
\begin{equation}\label{eq:large-two-y}
 Y\le \min\{6d\log(2M),\,ML/P^2\}.
\end{equation}
The second bound uses at most one prime per square multiplier
in each terminal interval and \(\log p\le L/2\).

At a nonexceptional equal-height prime its entire norm
power divides an integer in \([a+1,a+2d-1]\).
The product of these powers divides \(\binom{a+2d-1}{2d-1}\)
because the primes exceed \(2d-1\).
Writing \(A_0=a/d\), its logarithm is less than
\(2d(1+\log(A_0+2)-\log2)\).
At \(p\le P\) apply the affine estimate
\eqref{eq:norm-affine}. Comparing with
\(\log|F(0)|>H\log(aN)\), and setting
\[
 \Delta=\log(aN/H)-\frac2h(1+\log(A_0+2)),
 \qquad \Delta_*=\Delta+\frac2h\log2 ,
\]
gives
\begin{equation}\label{eq:large-norm}
 \Delta_*<\log(P/H)+\frac{\pi(P)L}{H}+\frac{Y}{H}+\frac{L}{4B}.
\end{equation}
Here \(\log P>8\), so \(S(P)<\log P\), and
\(2ML/(HP^3)\le L/(4B)\). The same estimate holds for \(\Delta\).

In \eqref{eq:large-parameter-range} we have
\(A_0+1>e^{U/2}/\sqrt{2U+2}>308\).
For \(h=2\), \(\Delta-U=\log(A_0/(A_0+2))-1-\log2>-1.70\),
and \(\Delta_*>U-1.01\).
For \(h\ge4\), \(\Delta-U\) attains its minimum on
\(4\le h\le 2.24(A_0+1)/(U+.47)\) at an endpoint, since
its derivative has the sign of \(2(1+\log(A_0+2))-h\).
At the lower endpoint the value exceeds \(.97\).
At the upper endpoint it is
\[
 \log\frac{U+.47}{2.24}+\log\frac{A_0}{A_0+1}
 -\frac{2(U+.47)}{2.24}\frac{1+\log(A_0+2)}{A_0+1}.
\]
This increases with \(A_0\). Substitute the lower bound
\(A_0+1>e^{U/2}/\sqrt{2U+2}\). The result exceeds \(1.6\)
at \(U=15\) and increases thereafter. Indeed, if
\(T=e^{U/2}/\sqrt{2U+2}\), then \(.46<T'/T<.5\).
The logarithmic derivative of
\((U+.47)(1+\log(T+1))/T\) is at most
\(1/(U+.47)+.5/(1+\log(T+1))-.46<0\).
In particular \(\Delta>U+.95\) for \(h\ge4\).
Now \(\pi(P)L/H<10\) by Lemma~\ref{lem:prime-estimates},
since \(P/H\le2.5\) and \(\log P>L/4+1.2762\).
By \eqref{eq:large-two-y}, for \(h\ge4\),
\[
 h<\frac{6(3.24U+1)}
           {U+.95-\log2.5-10-3.24U/8192}<61 .
\]
The fraction decreases from \(U=15\). Thus \(h\le60\).

We next refine the digit loss in \eqref{eq:pade-norm}.
For \(p>P'\ge2H+d\) and a positive lifted valuation \(e\),
one has \(e\ge1/H\) and \(e-\mu\ge1/(2H)\).
The progression in \eqref{eq:pade-progression} has valuation
at most \(N\mu+\lfloor\log_p(bm+j)\rfloor\), as \(p\nmid d\).
The floor first increases beyond \(L_p=\lfloor\log_p M\rfloor\)
only at \(j-1\ge b>2H\). Later increases are more than \(2H\)
steps apart. The penalty \((j-1)(e-\mu)\) absorbs all these
increases. Thus above \(P'\) the per-prime digit loss in the
averaged bound is at most \(L_p\log p/N\), with \(L_p\le3\).
In particular the digit loss is at most
\begin{equation}\label{eq:large-digit}
 \frac{\pi(P')V}{N}+\frac3N\sum_{p>P'}v_p(F(0))\log p .
\end{equation}

Put \(\delta=H/N\) and take \(P'=(2+2/h)H=rH\).
For nonexceptional primes above \(P'\),
Lemma~\ref{lem:two-heights} puts their entire norm product
in the binomial coefficient
\(\binom{b+H+d-1}{2H+2d-1}\).
With \(c=b/H>6.9\), its logarithm divided by \(H\) is at most
\(r(1+\log((c+r/2)/r))\). This increases with \(r\le3\) and is
less than \(3(t-.595)\), since \(c=e^t/h\). For the primes
below \(P'\), the first term of \eqref{eq:large-digit} is
at most \(12\delta\) using \(\log P'>L/4+2.59\).
The exceptional primes contribute at most \(6\delta L/(Br^3)\).
Therefore the digit loss is less than
\(\delta(9t+J_0(U))\), where \(J_0(U)=6.65+6(3.24U)/16384\).

When \(d\ge50\) the weaker angular bound gives
\(\log\Lambda>U+.98t-1.64\).
Use \(t>(U-\log(2U+2))/2\), \(g_0/d<.025\), and
\((2/\pi)\arcsin(e^{-t/2})g_0<.19\).
The last bound follows by substituting \(g_0<1.266+.256U\)
and the lower bound on \(t\) into
\(\arcsin x\le x/\sqrt{1-x^2}\). The resulting expression
decreases from \(U=15\). Split the prime sum at \(HV>e^8\).
If \(\delta\le .045\), the gain exceeds the two upper bounds
by more than
\[
 (.98-9\delta)t-\log\delta-\log(3.24U)-2.101-J_0(U)\delta.
\]
Substituting \(\delta=.045\) and \(t=(U-\log(2U+2))/2\)
decreases the expression. The resulting bound is above \(.10\)
at \(U=15\) and increases thereafter.
For \(d<50\), \(H<3000\), so the digit bound is below \(.14\).
The gain exceeds \(U+t/3-.71\), and \(t>U-\log50\).
It exceeds \(\log(3000\cdot3.24U)+.461+.14\) at \(U=15\) and
increases faster. Therefore
\begin{equation}\label{eq:large-degree}
 h\le60,\qquad \delta>.045,\qquad \Delta_*>U-1.01 .
\end{equation}

\subsection*{A count of terminal square multiples}

We use a sharper count in the range \eqref{eq:large-parameter-range}.
For either \(i=N-1,N\), set \(X=bi\), and count primes \(p>P\)
with \(X-2H\le qp^2\le X\), \(q\) a positive integer.
Divide the range up to \(\sqrt X\) into intervals
\([D,\eta D)\) with \(D=\eta^j P\) and \(\eta=9/8\).
The integer points \((p,q)\) lie between \(X/p^2\) and
\(X/p^2-\epsilon\), where \(\epsilon=2H/D^2\).
Each prime permits at most one multiplier and each multiplier
at most one prime. The mean value theorem gives
\(1\le2X|p_2-p_1|/D^3+\epsilon\) for distinct points.
Since \(\epsilon\le1/8\), their abscissas have spacing at least
\(g_D=\max(2,7D^3/(16X))\).

Partition each interval further into blocks of width
\(\rho=(D^4/X)^{1/3}\), with number at most
\(.132(X/D)^{1/3}\), since \(\rho/D\le X^{-1/6}<.007\).
For three points in a block the integer area determinant
is the determinant for \(X/x^2\) plus that for the vertical
errors. The first is at most \(3/4\) in absolute value by
the second divided difference bound \(3X/D^4\), times
\(\rho^3/4\). The error determinant is at most
\(\epsilon\rho\le(2H/X)^{1/3}<.01\). Thus all points in a
block are collinear. For \(n\ge3\) points, select a middle
one whose horizontal distances to both extremes are at least
\((n-2)g_D/2\). The chord of \(X/x^2\) through the extreme
abscissas is at most \(\epsilon\) above the curve there.
The lower curvature \(6X/(\eta D)^4\) gives
\[
 n\le 2+2\eta^2\sqrt{2/3}\,
                 \frac{D\sqrt{H/X}}{g_D}.
\]
The same bound holds for smaller \(n\).

Sum these bounds with the weight \(\log(\eta D)\).
The constant term from each block sums to at most
\[
 7(X/P)^{1/3}(\log P+3.06),
\]
since \(.264/(1-\eta^{-1/3})<7\) and
\(\log\eta/(1-\eta^{-1/3})<3.06\).
For the second term split at \(D_*=(32X/7)^{1/3}\),
where \(g_D\) switches. The powers of \(D\) in the two
geometric sums are \(2/3\) and \(-7/3\).
Their combined unweighted coefficient is bounded using
\[
 .132\eta^2\sqrt{2/3}(32/7)^{2/9}
 \left(\frac1{1-\eta^{-2/3}}+\frac1{1-\eta^{-7/3}}\right)<3.4 .
\]
The first sum has weight at most \(\log(\eta D_*)\).
For the second, \(D^{-7/3}\log(\eta D)\) decreases above
\(D_*\). Shifting its start to \(D_*\) bounds the mean weight
by \(\log(\eta D_*)+(\log\eta)\eta^{-7/3}/(1-\eta^{-7/3})
<(\log X)/3+1\).
Summing over both orders and using \(X\le M\) yields
\begin{equation}\label{eq:large-square-charge}
 \frac{Y}{H}<
 \frac{14(\log P+3.06)}{(P/H)^{1/3}B^{1/3}}
 +\frac{6.8 M^{1/18}}{\sqrt H}(L/3+1).
\end{equation}

\subsection*{The final comparison}

Put \(v=N e^{-t}=sh/\delta>6/5\). Then
\[
 M=\frac{\delta N^3}{hv},\qquad B=hv\delta^3N .
\]
Using \(\log P>U+\log(.09)\), \(L<3.24U\),
and \(\Delta_*>U-1.01\) in \eqref{eq:large-norm}
gives \(Y/(HL)>.05\). Explicitly, a lower bound is
\[
 \frac{U-1.01-\log2.5-
       \dfrac{2.5(3.24U)}{U+\log(.09)}
          \left(1+\dfrac{1.2762}{U+\log(.09)}\right)
       -3.24U/8192}{3.24U}.
\]
This increases for \(U\ge15\) and exceeds \(.0502\) there.
By \eqref{eq:large-two-y},
\(Y/(HL)\le 1/(v h(2+1/h)^2\delta^2)\). Hence
\(\delta^2<4/3\), since \(h(2+1/h)^2\ge12.5\).
In particular \(\delta<1.4\) and \(L<3U\).

For \(\ell\le\delta\le u\) and \(v\ge w\), define
\[
\begin{aligned}
 V_0&=3U+\log(u/(2w))+.001,\\
 Z&=\frac1\ell\sum_{j=4}^8(u/(2w))^{1/j}e^{(3/j-1)U}
       +\frac{3U}{\ell\log2}(u/(2w))^{1/9}e^{-2U/3},\\
 C_0&=1.017\{5+\min(2.5,(2w\ell^2)^{-1/3})+Z\},\\
 A(U,\ell,u,w)&=U-1.01-\log2.5-C_0-V_0/8192 .
\end{aligned}
\]
At \(p\le P\), the extra term in the affine norm count is
\(\sum_{j\le L/\log2}\vartheta(\min(P,M^{1/j}))\), where
\(\vartheta(x)=\sum_{p\le x}\log p\).
By Rosser and Schoenfeld,
\(\vartheta(x)<1.017x\) for \(x>0\) and
\(\vartheta(x)>.98x\) for \(x\ge7481\)
\cite[Theorems 9 and 10]{rosser-schoenfeld}.
In the upper estimate, \(j=1,2\) contribute at most
\(1.017(2P/H)\) after division by \(H\), the term \(j=3\)
uses \(M^{1/3}/H\le(2w\ell^2)^{-1/3}\), and the remaining
terms contribute at most \(1.017Z\).
Thus \eqref{eq:large-norm} with this count gives \(Y/H>A\).
On the other hand, \eqref{eq:large-square-charge} gives
\(Y/H<W(U,\ell,w)\), with
\[
 W(U,x,w)=e^{-U/3}\left\{
 \frac{14(U+\log(2.5x)+3.06)}{(5w)^{1/3}x}
 +\frac{6.8(U+\tfrac13\log(x/(2w))+1)}{(2w)^{1/18}x^{4/9}}
 \right\}.
\]
In obtaining this formula one can first take \(x=\delta\).
The function \(W\) decreases in \(x\) on the ranges used
here, and \(xW\) increases. In particular \(Y/N<uW(U,u,w)\).
At fixed \(\ell,u,w\), \(W\) decreases in \(U\ge15\)
while \(A\) increases. These statements follow directly
by differentiating the logarithmic factors and powers.

At \((\ell,u,w)=(.21,1.4,6/5)\) the difference \(A-W\)
exceeds \(.074\) at \(U=15\), so \(\delta<.21\).
On the resulting range \([.045,.21]\) with \(w=6/5\),
\(A/V_0\) increases with \(U\). To check this, write
\(A/V_0=(U-c(U))/(3U+\beta)-1/8192\),
where \(c(U)=1.01+\log2.5+C_0\) decreases and
\(\beta=\log(.21/2.4)+.001\) satisfies \(\beta+3c(U)>0\).
The multiplier bound in \eqref{eq:large-two-y} now gives
\(sh<.675/\delta\), since \((2/25)V_0/A<.675\) at \(U=15\).
For larger \(\delta\), the table below uses the stronger
bound \(sh<uw\) when \(A>W\) excludes \(v\ge w\).

The weaker angular bound in Lemma~\ref{lem:angles} gives
\(\log\Lambda>U+t-1.60\). Here \((t+g_0)/d<.009\)
from \(d>.045N/60\), the angle loss is below \(.19\)
as before, and \(\log C>t-1.39\).
For each row \([\ell,u]\) put \(V_0=V_0(U,u,6/5)\) and set
\[
 Q=\min\{V_0,\ U+.127+W(U,\ell,6/5)+V_0/8192\}.
\]
Both entries bound the large-prime log norm divided by \(H\).
For the second entry, the equal-height prime contribution
is at most \((2/h)(U+.127)\), since
\(A_0+2<(5/6)N+1/6\), and the other contributions use
\eqref{eq:large-square-charge} and the cube bound.
As \(HQ>P\), splitting the prime sum at \(HQ\) bounds it
by \(\log(HQ)+.461\).

The digit estimate \eqref{eq:large-digit} has \(L_p\le2\)
above \(P\), except that \(L_p\le3\) when \(p\le M^{1/3}\).
The extra loss is at most
\(\max(0,\vartheta(M^{1/3})-\vartheta(P))/N\), which is
bounded by \(\max(0,1.017(u/2.4)^{1/3}-1.96\ell)\).
Thus the digit loss is at most the following, with
\(z=U+\log(2.5\ell)\):
\[
\begin{aligned}
 D_0(U,\ell,u)
  &=u\left\{\frac{2.5V_0}{z}\left(1+\frac{1.2762}{z}\right)
                       +2(U+.127)+V_0/4096\right\}\\
  &\qquad+2uW(U,u,6/5)+\max(0,1.017(u/2.4)^{1/3}-1.96\ell).
\end{aligned}
\]
For the prime count in this formula we use that
\(r/\log(rH)(1+1.2762/\log(rH))\) increases for
\(2\le r\le2.5\). Since \(H=she^t\), the surplus over
\eqref{eq:pade-norm} is greater than
\(U-\log(JQ)-2.061-D_0\) where \(J\) bounds \(sh\).
The following endpoint comparisons complete the proof:
\[
\begin{array}{c|c|r|c|c}
{}[\ell,u]&w&J&
\substack{A-W\\\text{at }U=15}&
\substack{U-\log(JQ)-2.061-D_0\\ \text{at }U=15}\\ \hline
{}[.045,.10]&\text{--}&15&\text{--}&>.36\\
{}[.10,.14]&9&1.26&>.109&>1.63\\
{}[.14,.18]&3.6&.648&>.117&>.83\\
{}[.18,.21]&1.8&.378&>.093&>.30
\end{array}
\]
The first row uses \(sh<.675/\delta\). In the other rows
the fourth column gives \(v<w\). The final column
increases with \(U\). Indeed the logarithmic growth rate of
either branch of \(Q\) is below \(.08\).
The rational factor \(V_0(1+1.2762/z)/z\) decreases with \(U\):
write it as \(3(1-k/z)(1+1.2762/z)\), where
\(k=\log(2.5\ell)-(\log(u/2.4)+.001)/3<.02\).
The \(W\) term decreases and the other terms of \(D_0\)
grow at rate at most \(.21(2+3/4096)<.43\).
Thus the positive surplus contradicts the Padé inequality.
The accompanying verifier reproduces the scalar comparisons
at the stated endpoints using rational enclosures for logarithms,
exponentials, and rational powers.
\end{proof}

